\documentclass[10pt,a4paper]{amsart}
\usepackage[margin=19mm]{geometry}
\usepackage{amsmath,amssymb,mathtools}
\usepackage[scr=rsfs,cal=euler]{mathalfa}
\usepackage{xfrac}
\usepackage[unicode,hidelinks,bookmarks=false]{hyperref}
\newcommand{\ie}{i.e.\ }
\newcommand{\cf}{cf.\ }
\newcommand{\resp}{respectively\ }
\newcommand{\Subsection}{\subsection}
\providecommand{\nobreakdash}{\nobreak\hbox{-}}
\setlength{\emergencystretch}{2em}
\multlinegap0pt
\usepackage{amssymb}
\usepackage{mathtools}
\usepackage[shortlabels]{enumitem}
\usepackage{xcolor}
\newtheorem{lemma}{Lemma}
\newtheorem{proposition}{Proposition}
\usepackage{cleveref}
\crefname{lemma}{Lemma}{Lemmas}
\crefname{proposition}{Proposition}{Propositions}
\newcommand{\Dir}{\operatorname{Dir}}
\newcommand{\Dn}{\mathcal D_n}
\newcommand{\Fq}{\mathbb F_q}
\newcommand{\HH}{\mathbb H}
\newcommand{\Mrd}{\mathcal M^{\mathrm{rd}}_{\HH_n}}
\newcommand{\PP}{\mathbb P}
\newcommand{\norm}[1]{\lVert #1\rVert}
\newcommand{\one}{\mathbf 1}
\title{Sharp refined-direction Kakeya estimates in finite Heisenberg groups}
\author{Thang Pham, Andrea Pinamonti, Dung The Tran,, Boqing Xue}
\date{Source: arXiv:2608.14059v1 (CC BY 4.0)}
\begin{document}
\maketitle

\noindent\emph{Source and licence.} Sharp refined-direction Kakeya estimates in finite Heisenberg groups. Version arXiv:2608.14059v1 (CC BY 4.0), distributed by arXiv under the Creative Commons Attribution 4.0 International licence, http://creativecommons.org/licenses/by/4.0/, as recorded in the arXiv licence metadata for this exact version and shown on the abstract page https://arxiv.org/abs/2608.14059v1. Full licence notice: \texttt{licenses/arxiv-2608.14059-CC-BY-4.0.txt}.\par\smallskip

\noindent\textit{Editorial scope.} Counted unit: the article's proposition prop:lower-bounds (source0.tex lines 832--838). The article's complete original proof of it is delivered with it (source0.tex lines 840--974, one \texttt{proof} environment of 135 lines). No mathematics is written, repaired or re-derived: the statement and the proof are the article's own lines, in the article's own notation, and no retained line is substituted or re-flowed.  Retained uncounted as the source-local context the proof needs: lines 236--242; lines 381--386; lines 664--666.  Nothing was omitted from any retained span, so the package carries no omission record.  No retained line contains a citation, so the wrapper carries no bibliography and no bibliography entry.  No retained line contains a figure environment, an image-inclusion command or drawing code, so no image is delivered and the package declares no asset dependency.  Editorial formatting only, in addition: an AMS \texttt{amsart} preamble that restates the article's own declarations and macros used by the retained lines, namely the environments \texttt{lemma}, \texttt{proposition} on separate counters, together with the standard packages those lines need.  The retained environments print in this document as Lemma 1, Proposition 1 in order of appearance, whereas the article prints the same items with the numbering of its own full document.  The complete native source of the article as distributed in the arXiv e-print for this exact version is bundled unchanged as \texttt{sources/207631/source0.tex}.
\begin{equation}\label{eq:cardinalities}
 \begin{aligned}
 |\HH_n(\Fq)|&=q^{r+1},\\
 |\PP^{r-1}(\Fq)|&=1+q+\cdots+q^{r-1},\\
 |\Dn|&=q+q^2+\cdots+q^r.
 \end{aligned}
\end{equation}

\begin{equation}\label{eq:def-critical-exponent}
 \norm{\Mrd F}_{\ell^v(\Dn)}
 \leq
 Cq^a
 \norm{F}_{\ell^u(\HH_n(\Fq))}
\end{equation}

\begin{lemma}\label{lem:line-geometry}
Every point of $\HH_n(\Fq)$ lies on exactly $|\PP^{r-1}(\Fq)|$ horizontal lines, and these lines have pairwise distinct refined directions. Two distinct points lie on at most one common horizontal line.
\end{lemma}

\begin{proposition}\label{prop:lower-bounds}
For every $(\alpha,\beta)\in[0,1]^2$,
\begin{equation}\label{eq:lower-bounds}
 \widetilde A_n^{\mathrm{rd}}(\alpha,\beta)
 \geq\Phi_n(\alpha,\beta).
\end{equation}
\end{proposition}

\begin{proof}
Fix $u=\frac{1}{\alpha}$ and $v=\frac{1}{\beta}$, with the usual convention when $\alpha=0$ or $\beta=0$. If \eqref{eq:def-critical-exponent} holds with the power $q^a$, then every nonzero test function $F$ satisfies
\begin{equation}\label{eq:test-function-principle}
 q^a
 \gtrsim_{n,u,v,a}
 \frac{\norm{\Mrd F}_{\ell^v(\Dn)}}
      {\norm{F}_{\ell^u(\HH_n(\Fq))}}.
\end{equation}
We give four test functions and estimate the corresponding quotients from below.

\emph{A point mass.}
Fix $\mathbf{p}_0=(\mathbf{z}_0,t_0)\in\HH_n(\Fq)$ and set
$F=\one_{\{\mathbf{p}_0\}}$.  By \cref{lem:line-geometry},
$\mathbf{p}_0$ lies on exactly
\[
 |\PP^{r-1}(\Fq)|=1+q+\cdots+q^{r-1} \geq q^{r-1}
\]
horizontal lines with pairwise distinct refined directions. On each of these directions, $\Mrd F=1$. Since $\norm{F}_{\ell^u(\HH_n(\Fq))}=1$, it follows that
\[
 \frac{\norm{\Mrd F}_{\ell^v(\Dn)}}
      {\norm{F}_{\ell^u(\HH_n(\Fq))}}
 \geq |\PP^{r-1}(\Fq)|^{\frac{1}{v}}
 \geq q^{\frac{r-1}{v}}.
\]
Thus, every admissible power satisfies $a\geq\phi_1(\alpha,\beta)$.

\emph{One horizontal line.}
Let $L$ be a horizontal line and take $F=\one_L$. The standard
parametrization of $L$ is injective in the field parameter, so $|L|=q$
and
\[
 \norm{F}_{\ell^u(\HH_n(\Fq))}=q^{\frac{1}{u}}.
\]
At $\vartheta_0=\Dir(L)$, the line $L$ itself is admissible in the
definition of the maximal operator. Thus,
\[
 \Mrd F(\vartheta_0)
 \geq\sum_{\mathbf{p}\in L}\one_L(\mathbf{p})
 =q,
\]
and every $\ell^v$-norm dominates this single coordinate. Hence,
\[
 \frac{\norm{\Mrd F}_{\ell^v(\Dn)}}
      {\norm{F}_{\ell^u(\HH_n(\Fq))}}
 \geq q^{1-\frac{1}{u}}.
\]
Hence, $a\geq\phi_2(\alpha,\beta)$.



\emph{A small set realizing every refined direction.}
We seek a set of $O_n(q)$ points in the central slice $V\times\{0\}$ that meets a horizontal line in every refined direction. A point $(z,0)$ lies on a horizontal line of refined direction $[w:c]$ precisely when $\sigma(z, w)=c$. Thus, it is enough to choose a small subset of $V$ that meets every affine hyperplane arising from such an equation.

Fix a basis $e_1,\ldots,e_r$ of $V$ and define
\begin{equation}\label{eq:axes-set}
 \mathcal B
 :=
 \bigcup_{j=1}^r
 \{s e_j:s\in\Fq\}.
\end{equation}
Then $|\mathcal B|=1+r(q-1) \leq r q$, and $\mathcal B$ meets every affine hyperplane in $V$. Indeed, if $\varphi:V\to\Fq$ is a nonzero linear functional and $\Pi=\{z:\varphi(z)=c\}$, choose $j$ with
$\varphi(e_j)\neq 0$; then
$\frac{c}{\varphi(e_j)} \cdot e_j\in\Pi\cap\mathcal B$.

Let $E:=\{(z,0): z\in\mathcal B\}$ and $F:=\one_E$. 
Fix an arbitrary $\vartheta=[w:c]\in\Dn$. Since
$w\neq0$ and $\sigma$ is nondegenerate, the functional
\[
 z\longmapsto\sigma(z,w)
\]
is nonzero. Consequently,
\[
 \Pi_\vartheta
 :=\{z\in V:\sigma(z, w)=c\}
\]
is an affine hyperplane. This hyperplane is independent of the chosen representative of the projective point $[w:c]$. Choose $z_\vartheta\in\mathcal B\cap\Pi_\vartheta$. Then
\begin{equation}\label{eq:witnessing-line-final}
 L_{(z_\vartheta,0), w}
 =
 \{(z_\vartheta+sw,sc):s\in\Fq\}
\end{equation}
has refined direction $
 [w:\sigma(z_\vartheta,w)]
 =[w:c]
 =\vartheta$, 
and it contains $(z_\vartheta,0)\in E$. Consequently,
\[
 \Mrd F(\vartheta)\geq1
 \qquad(\vartheta\in\Dn),
\]
which together with $|E|=|\mathcal{B}| \leq r q$ implies that
\[
 \frac{\norm{\Mrd F}_{\ell^v(\Dn)}}
      {\norm{F}_{\ell^u(\HH_n(\Fq))}}
      \geq
 \frac{|\Dn|^{\frac{1}{v}}}{|E|^{\frac{1}{u}}}
 \geq
 r^{-\frac{1}{u}}q^{\frac{r}{v}-\frac{1}{u}}
 \gtrsim_n q^{\frac{r}{v}-\frac{1}{u}}.
\]
Thus, $a\geq\phi_3(\alpha,\beta)$.


\emph{The constant function.}
Finally, take $F\equiv1$ on $\HH_n(\Fq)$. Since
$|\HH_n(\Fq)|=q^{r+1}$,
\[
 \norm{F}_{\ell^u(\HH_n(\Fq))}
 =q^{\frac{r+1}{u}}.
\]
Every refined direction is realized by a horizontal line, and every
horizontal line contains $q$ points. Hence,
\[
 \Mrd F(\vartheta)=q
 \qquad(\vartheta\in\Dn).
\]
Therefore, by the third identity in \eqref{eq:cardinalities},
\[
 \frac{\norm{\Mrd F}_{\ell^v(\Dn)}}
      {\norm{F}_{\ell^u(\HH_n(\Fq))}}
 =
 \frac{q|\Dn|^{\frac{1}{v}}}{q^{\frac{r+1}{u}}}
 \geq
 q^{1+\frac{r}{v}-\frac{r+1}{u}}
 =
 q^{\phi_4(\alpha,\beta)}.
\]
Hence, $a\geq\phi_4(\alpha,\beta)$. Every admissible $a$ therefore
dominates all four affine functions, so
\[
 a\geq\max_{1\leq j\leq4}\phi_j(\alpha,\beta)
 =\Phi_n(\alpha,\beta).
\]
Taking the infimum over admissible powers proves the proposition.
\end{proof}

\end{document}
